% Options for packages loaded elsewhere
\PassOptionsToPackage{unicode}{hyperref}
\PassOptionsToPackage{hyphens}{url}
%
\documentclass[
  12pt,
]{article}
\usepackage{amsmath,amssymb}
\usepackage{lmodern}
\usepackage{iftex}
\ifPDFTeX
  \usepackage[T1]{fontenc}
  \usepackage[utf8]{inputenc}
  \usepackage{textcomp} % provide euro and other symbols
\else % if luatex or xetex
  \usepackage{unicode-math}
  \defaultfontfeatures{Scale=MatchLowercase}
  \defaultfontfeatures[\rmfamily]{Ligatures=TeX,Scale=1}
\fi
% Use upquote if available, for straight quotes in verbatim environments
\IfFileExists{upquote.sty}{\usepackage{upquote}}{}
\IfFileExists{microtype.sty}{% use microtype if available
  \usepackage[]{microtype}
  \UseMicrotypeSet[protrusion]{basicmath} % disable protrusion for tt fonts
}{}
\makeatletter
\@ifundefined{KOMAClassName}{% if non-KOMA class
  \IfFileExists{parskip.sty}{%
    \usepackage{parskip}
  }{% else
    \setlength{\parindent}{0pt}
    \setlength{\parskip}{6pt plus 2pt minus 1pt}}
}{% if KOMA class
  \KOMAoptions{parskip=half}}
\makeatother
\usepackage{xcolor}
\IfFileExists{xurl.sty}{\usepackage{xurl}}{} % add URL line breaks if available
\IfFileExists{bookmark.sty}{\usepackage{bookmark}}{\usepackage{hyperref}}
\hypersetup{
  hidelinks,
  pdfcreator={LaTeX via pandoc}}
\urlstyle{same} % disable monospaced font for URLs
\usepackage[margin=1in]{geometry}
\setlength{\emergencystretch}{3em} % prevent overfull lines
\providecommand{\tightlist}{%
  \setlength{\itemsep}{0pt}\setlength{\parskip}{0pt}}
\setcounter{secnumdepth}{-\maxdimen} % remove section numbering
\ifLuaTeX
  \usepackage{selnolig}  % disable illegal ligatures
\fi

\author{}
\date{}

\begin{document}

\hypertarget{personal-statement-nus-young-fellowship-programme-2026}{%
\section{Personal Statement — NUS Young Fellowship Programme
2026}\label{personal-statement-nus-young-fellowship-programme-2026}}

\textbf{Topic 2: A current social issue of significant personal
relevance}

\begin{center}\rule{0.5\linewidth}{0.5pt}\end{center}

In late 2023, I set foot inside the DENSO factory as part of the DENSO
Factory Hacks — a collaborative program between FPT and DENSO. I had
visited a handful of factories before, but what I encountered at DENSO
was something fundamentally different: cutting-edge technology,
optimized processes, and above all — humans working in harmony with
machines, supported by technologies I had never seen deployed at such
scale: augmented reality headsets for worker training, autonomous guided
vehicles navigating the floor, and intelligent production management
systems. What struck me most deeply was not the precision of robotic
arms or the seamless assembly lines, but DENSO’s operating philosophy:
\emph{“People are at the center.”} Here, technology did not exist to
eliminate humans, but to amplify their capabilities — and that
philosophy has since become the compass guiding every research direction
I have pursued.

Yet when I stepped back from DENSO’s factory floor and looked across the
region, a troubling question surfaced: if this was the pinnacle that
global corporations had reached, where did the tens of thousands of
small and medium factories across Southeast Asia — where workers still
operated largely through intuitive experience and manual processes —
stand on this journey? That question has driven me ever since. Every
project I have undertaken — from an intelligent fault-query system to a
multi-agent vehicle coordination framework — traces its roots back to
that experience at DENSO and the aspiration to bridge that gap.

\hypertarget{the-digital-transformation-gap-in-southeast-asian-manufacturing}{%
\subsection{The Digital Transformation Gap in Southeast Asian
Manufacturing}\label{the-digital-transformation-gap-in-southeast-asian-manufacturing}}

Southeast Asia is one of the world’s fastest-growing manufacturing hubs.
In Vietnam alone, industrial production constitutes a substantial share
of GDP and employs millions of workers. Yet a thought-provoking paradox
persists: while cutting-edge AI systems can now autonomously coordinate
robot fleets, predict equipment failures before they occur, and optimize
energy consumption in real time, the vast majority of factories in the
region still depend on manual workflows, paper-based documentation, and
reactive maintenance.

This gap is not merely technological. It is \textbf{systemic} —
encompassing a critical shortage of engineers capable of translating AI
research into practical industrial solutions, an absence of locally
contextualized models, and a cultural rift between academic communities
and shop-floor practitioners. But perhaps the most pressing dimension
runs in the opposite direction: as modern technology advances,
\textbf{millions of workers} whose labor has long been Southeast Asia’s
competitive advantage face the risk of displacement if digital
transformation proceeds through substitution rather than collaboration.

Growing up in Vietnam and studying Mechatronics Engineering at Hanoi
University of Science and Technology (HUST), I have witnessed this
disconnect daily — not through textbooks, but through factory visits and
conversations with workers and field engineers on the ground.

\hypertarget{my-journey-from-sensors-to-intelligent-systems}{%
\subsection{My Journey: From Sensors to Intelligent
Systems}\label{my-journey-from-sensors-to-intelligent-systems}}

Throughout my undergraduate years, every project I pursued has followed
a consistent direction: \emph{harnessing cutting-edge, emerging
technologies to solve real industrial problems — enhancing operational
efficiency and automating manufacturing processes.}

\textbf{UWB Warehouse System (2022):} It began when I developed a
warehouse management system using Ultra-Wideband (UWB) technology to
deliver real-time indoor positioning for asset tracking. The project
earned HUST’s Innovation and Scientific Research Award and, more
importantly, taught me my first lesson about the complexity of
integrating hardware sensors with cloud data architectures in real-world
environments.

\textbf{DENSO Factory Hacks (late 2023 – early 2024) — Turning Point:}
The turning point came during the DENSO Factory Hacks, where my team and
I developed an intelligent query system for maintenance engineers. DENSO
already operated an internal system that logged machine faults — but
engineers still spent considerable time searching and cross-referencing
errors within the vast data repository. Our solution allowed them to ask
questions in natural language and receive precise answers from the
entire historical fault database — compressing hours of manual lookup
into seconds. The system earned Third Prize, but the real reward was
first-hand exposure to factory-floor operations: it revealed both the
enormous potential of AI in manufacturing and the formidable barriers to
its real-world deployment.

\textbf{Multi-Agent RL for AGV Coordination (2025) — Research Frontier:}
That experience directly fueled my most ambitious research endeavor:
designing a \textbf{multi-agent reinforcement learning (MARL) framework}
for autonomous guided vehicle (AGV) coordination in smart factory
environments. The architecture combines \textbf{Proximal Policy
Optimization (PPO)} with \textbf{Graph Neural Networks (GNN)} over a
heterogeneous super-graph representation, enabling multiple AGVs to
simultaneously learn cooperative behaviors — task allocation, collision
avoidance, and deadlock resolution — through training rather than
explicit programming. Notably, integrating classical PID controllers
with modern RL agents demonstrated that legacy control and AI are
\textbf{complementary}, opening a viable pathway for factories seeking
an incremental transition to intelligent systems without rebuilding from
scratch.

\hypertarget{why-this-matters-a-vision-for-human-ai-collaborative-manufacturing}{%
\subsection{Why This Matters: A Vision for Human-AI Collaborative
Manufacturing}\label{why-this-matters-a-vision-for-human-ai-collaborative-manufacturing}}

The projects I have undertaken share a unifying philosophy: the fusion
of advanced technology and human capability is the key to the
sustainable development of industry.

This is not merely idealism — it is a survival imperative for the
region. Southeast Asia has long competed globally through an abundant,
cost-effective workforce. But as automation proliferates, that advantage
becomes a vulnerability if digital transformation pursues replacement
rather than collaboration. I believe the answer lies in building
\textbf{collaborative systems} where AI elevates human capability —
transforming shop-floor workers into intelligent operators, and
converting intuitive experience into structured, decision-ready data.

This demands a fundamental shift in thinking. Rather than asking
\emph{“Where can machines replace humans?”}, we must ask \emph{“What do
humans need to perform even better when machines become their
teammates?”} The intelligent query system I built at DENSO did not
eliminate maintenance engineers — it transformed them into
decision-makers empowered by AI. Similarly, the multi-agent AGV
framework does not remove humans from the supervisory loop — it
optimizes task allocation and movement coordination at a scale exceeding
human cognitive capacity, freeing supervisors from reactive
micro-decisions to focus on strategic oversight and exception handling —
domains where human judgment remains irreplaceable.

My vision for Industry 4.0 in developing economies rests on three
pillars:

\begin{enumerate}
\def\labelenumi{\arabic{enumi}.}
\item
  \textbf{Collaboration over replacement}: The true value of AI in
  manufacturing lies in creating a complementary intelligence layer —
  where the practical experience of operators and the data processing
  capability of AI resonate together. When maintenance engineers query
  the system in natural language, or supervisors receive optimal
  allocation recommendations from RL algorithms — that is two forms of
  intelligence complementing each other to make decisions more accurate
  than either could achieve alone.
\item
  \textbf{Transform incrementally, not disruptively}: Developing nations
  cannot — and should not — wholesale import automation models from
  advanced economies. Solutions must accommodate limited infrastructure
  and varying digital readiness, where each step of transformation is
  accompanied by reskilling and upskilling the existing workforce,
  rather than mass layoffs.
\item
  \textbf{Bridging research and real-world deployment}: Most AI research
  in manufacturing exists only on paper. The region needs engineers who
  command both the theoretical foundations of AI and understand
  real-world operational challenges — and who can translate academic
  knowledge into deployable industrial solutions. That is precisely the
  profile I am building for myself.
\end{enumerate}

\hypertarget{why-nus-young-fellowship-programme}{%
\subsection{Why NUS Young Fellowship
Programme}\label{why-nus-young-fellowship-programme}}

The NUS Young Fellowship Programme, with its focus on “AI and Graduate
Research,” represents a rare opportunity to accelerate this vision from
idea to impact. The College of Design and Engineering (CDE) at NUS
houses world-leading research groups in robotics, intelligent systems,
and advanced manufacturing — domains directly aligned with my research
trajectory. I am particularly eager to explore the work of
\textbf{Prof.~ZHAO Lin} on optimization and control of multi-agent
systems, \textbf{Prof.~LOW Kian Hsiang} on machine learning for
multi-agent robotic systems, and the \textbf{Control, Intelligent
Systems \& Robotics (CISR)} group — research directions intimately
connected to the AGV coordination challenge I am investigating.

Singapore itself — with its national Advanced Manufacturing strategy and
position at the heart of Southeast Asia — offers a living laboratory
where I can observe how a small but highly advanced nation has
successfully transformed its manufacturing sector. My ambition extends
beyond learning: I intend to bring methodologies, frameworks, and
networks from NUS back to Vietnam — where thousands of enterprises stand
at the threshold of their digital transformation journey.

\hypertarget{long-term-vision}{%
\subsection{Long-Term Vision}\label{long-term-vision}}

Through this programme, I seek not only to deepen my research
capabilities but also to build the academic foundation and international
network necessary to pursue graduate studies. My ultimate goal is to
become a \textbf{bridge-builder} between cutting-edge AI research and
industrial practice — helping Southeast Asian factories not merely catch
up, but \textbf{leapfrog} into Industry 4.0, in a way that is
\textbf{inclusive}: leaving no one behind.

\end{document}
